\documentclass[12pt,a4paper]{article}
\usepackage{graphicx} % Required for inserting images
\usepackage{chemfig}
\usepackage{lewis}
\author{Svein-Tore Narvestad}
\date{September 2025}

\begin{document}
\textbf{\large{Løsningsforslag til prøve om kjemiske bindinger:}}\\\\\\
1.	Hva menes med begrepet elektronegativitet og hvor finner vi grunnstoffene som har høy/lav elektronegativitet?\\\\
Elektronegativitet er et mål på grunnstoffets evne til å trekke på elektroner i en kjemisk binding.  Elektronegativiteten øker utover i periodene og minker nedover i gruppene.  Grunnen til at det øker utover i periodene er at grunnstoffene da får flere elektroner i ytterste skall.  De nærmer seg da å få oppfyllt åtteregelen, og de holder derfor godt på elektronene sine.  Grunnen til at elektronegativiteten minker nedover i gruppene er at de da får et nytT elektronskall for hver ny periode.  Da blir avstanden til kjernen større, og dermed den elektrostatiske tiltrekkingskraften mindre.\\\\
2.	Hva menes med ioniseringsenergi? Hvilke grunnstoff har høy ioniseringsenergi?\\\\
Vi vet at atomene har en evne til å trekke på elektroner.  Det betyr at det  kreves en viss energi for å fjerne et elektron fra et atom. Og den energien som trengs for å fjerne et atom kalles for ioniseringsenergien.  Enkelt forkLart er dette energien som trengs for at atomet skal bli et ion.  Så er det en sammenheng melLom ioniseringsenerig og elektronegativiteten.  Atomer med høy elektronegtivitet har også høy ioniseringsenergi.  Så har vi også begreper som  1. ioniseringsenergi, og 2. ioniseringsenergi.  1. ioniseringsenergi er energien for å fjerne det første elektronet, og 2. ioniseringsenergi er energien for å fjerne det andre elektronet.  Disse kan variere mye.  Tenk på Na.  Det har ett elektron i ytterste skall. Det er lett å fjerne det første elektronet, og derfor er 1. ioniseringsenergi lav (495,8kJ/mol). Når Na har mistet ett elektron har det oppfylt åtteregelen og derfor er det vanskelig å fjerne et elektron til. Dette ser vi på verdien til 2. ioniseringsenergi.  Den er lik 4562kJ/mol.\newpage
3. Forklar hvorfor vann er en dipol.\\\\
For å forklare dette må vi se på vannets molekylstruktur:\\\\
\begin{center}
    \chemfig{H^{\delta+}-[:37.75]\charge{45=\:,135=\:,100:12pt=$2\delta^-$}{O}-[:-37.75]H^{\delta+}} 
\end{center}\vspace{1cm}
Skal et molekyl som består av 3 atomer ikke være en dipol, må atomene ligge på en rett linje.  Men det gjør det ikke.  Og grunnen til det er at oksygenatomet i vannmolekylet har 4 elektroner (2 elektronpar) som ikke deltar i binding.  Disse vil prøve å stå så langt fra hverandre som mulig (fordi de har lik ladning, de frastøtes hverandre).   Dette fører til at H atomene får en liten positiv ladning, og oksygenet en liten negativ ladning.  Molekylet blir dermed en dipol.\\\\
4.	Hva menes med hydrogenbinding?  Nevn 3 stoff som har hydrogenbinding.\\\\
Molekyler med mer enn 2 atomer der sentralatomet (det midtersete atomet) har elektroner i ytterste skall som ikke deltar i binding er dipoler.  Og da har man bindinger mellom mellom molekylene.  SlikE bindinger kaller man for dipol-dipol binding.  For noen slike molekyler, som inneholder hydrogen og ett av atomene N,O og F så viser det seg at denne dipol-dipol bindingen er spesiell sterk.  Og denne bindingen har derfor fått et eget navn, nemlig hydrogenbinding.  Tre molekyler som har hydrogenbinding: $H_2O, NH_3$ og $HF$.        \\\\
5.	 Gitt følgende stoffer: $H_2O, NH_3, CO_2,CH_4$ og NaCl.
Hva er elektronegativitetsverdien for grunnstoffene i stoffene over?\\\\
Grunnstoffene i molekylene/saltne over er  H,N, O, C, Na og Cl. Disse grunnstoffenE har følgende verdier for elektrronegativiteten:  E(H)=2,1  E(N)=3,5 E(Cl)=3,0 E(C)=2,5 E(Na)=0,9 E(O)=3,5 \\\\\newpage
6.	Hvilken bindingstype har vi mellom molekylene i stoffene $H_2O, NH_3, CO_2,CH_4$ og NaCl?\\\\
$\Delta E(O,H)=3,5-2,1=1,4$  Siden $0,4<\Delta E<1,7.$ så har vi polar kovalent binding \\
$\Delta E(N,H)=3,0-2,1=0,9$  Siden $0,4<\Delta E<1,7.$ så har vi polar kovalent binding\\
$\Delta E(O,C)=3,5-2,5=1,0$  Siden $0,4<\Delta E<1,7.$ så har vi polar kovalent binding\\
$\Delta E(C,H)=2,5-2,1=0,4$  Denne er litt vanskelig.  For når $\Delta E=0,4$ så er den hverken større eller mindre enn 0,4.  Dette er et eksempel på en svak polar binding.  Men her vil jeg gi riktig svar på både ren kovalent binding og polar kovalent binding.  Ionebinding blir helt feil. \\
$\Delta E(Cl,Na)=3,0-0,9=2,1$ Siden $\Delta E>2$ så er dette en ionebinding.\\\\
Her står tegnet $\Delta$ for forskjell. Så $\Delta E$ står for forskjell i elektronegativitet.
\\\\
7.	Tegn elektronprikkstrukturen for stoffene $H_2O, NH_3, CO_2, CH_4$ og NaCl.\\\\ 
Når man skal tegne elktronprikkstruktruren for et molekyl så skal vi tenge atomet med tildels riktige vinker, og vi må også elektroner som ikke deltar i binding for sentralatomet.  dette blir viktig i neste oppgave.\\\\\\
$H_2O$:\\\\
\begin{center}
    \chemfig{H-[:37.75]\charge{45=\:,135=\:}{O}-[:-37.75]H} 
\end{center}\vspace{1cm}
$NH_3$:
\begin{center}
    \chemfig{H-\charge{90=\:}{N}(-[:-90]H)(-[:0]H)} 
\end{center}\newpage
$CO_2$:
\begin{center}
    \chemfig{\lewis{O}{}{}{.}{.}{}{}{.}{.}=C=\lewis{O}{}{}{.}{.}{}{}{.}{.}}
\end{center}\vspace{1cm}
$CH_4$:
\begin{center}
    \chemfig{C(-[:-90]H)(-[:0]H)(-[:90]H)(-[:180]H)}
\end{center}\vspace{1cm}
NaCl:
\begin{center}
    $\Biggr[Na \Biggr]^+$$\Biggr[\lewis{Cl}{.}{.}{.}{.}{.}{.}{.}{.} \Biggr]^-$
\end{center}\vspace{1cm}
8.	Er noen av molekylene dipoler? Begrunn svaret.\\\\
Alle molekyølene over består av mer enn 3 atomer (NaCl er ikke et atom).  Skal slike mkolekyler være endipol må de ha en lasningsforkyvning i sttrukturen.  og det oppnår de dersom sentralatomet til molekylet har elektroner som ikke deltar i binding. For moleklylene over er det $H_2O$ og $NH_3$ som er dipoler.  
\end{document} 
